% topic_05.tex — Content only. Compiled via main.tex.
% ── No \documentclass, \begin{document} or \end{document} here ──




\chapteropener{5}{Work, Energy and Power}{%
  \item Work Done by a Force
  \item Conservation of Energy and Efficiency
  \item Power and \texorpdfstring{$P = Fv$}{P = Fv}
  \item Gravitational Potential Energy
  \item Kinetic Energy
  \item Energy Transfers and Problem Solving
}

% ===== SECTION 1: WORK =====
\section{Work Done by a Force}

\begin{definitionbox}{Work Done}
\textbf{Work done} is defined as the product of the force and the displacement \textbf{in the direction of the force}:
$$W = Fs\cos\theta$$
\begin{itemize}[itemsep=2pt]
  \item $F$: applied force (N); \quad $s$: displacement (m); \quad $\theta$: angle between force and displacement.
  \item Unit: joule (J) $\equiv$ N m.
  \item Work is a \textbf{scalar} quantity.
  \item If force and displacement are parallel ($\theta = 0$): $W = Fs$.
  \item If force is perpendicular to displacement ($\theta = 90^\circ$): $W = 0$ (no work done).
\end{itemize}
\end{definitionbox}

\begin{center}
\begin{tikzpicture}[scale=1.1, every node/.style={font=\small}]
  % Object
  \fill[midblue!20, draw=darkblue, thick, rounded corners=3pt]
    (0,-0.3) rectangle (0.8, 0.3);
  % Displacement arrow
  \draw[-latex, darkblue, very thick] (0.8,-0.5) -- (4.0,-0.5)
    node[anchor=north, midway] {$s$};
  % Force arrow at angle theta
  \draw[-latex, pinkframe, very thick] (0.4,0) -- (2.6, 1.3)
    node[anchor=south west] {$F$};
  % Horizontal component dashed
  \draw[-latex, , thick, darkgray] (0.4,0) -- (2.6,0);
  \draw[dashed, darkgray] (2.6,0) -- (2.6,1.3);
  % Angle arc
  \draw[accentgold, thick] (1.0,0) arc (0:30:0.6);
  \node[accentgold] at (1.55, 0.18) {$\theta$};
  % Component label
  \node[darkgray, font=\footnotesize] at (1.5,-0.25) {$F\cos\theta$};
\end{tikzpicture}
\end{center}

\begin{keybox}{Work Done Against Gravity}
When an object of mass $m$ is lifted vertically through height $\Delta h$:
$$W = mg\Delta h$$
The force required equals the weight $mg$, and the displacement is $\Delta h$ in the same direction — so $\theta = 0$ and $\cos\theta = 1$.
\end{keybox}

\begin{warningbox}{Common Mistake}
Work is done by the \emph{component} of force in the direction of motion — not the full force. A force perpendicular to motion (e.g.\ the normal contact force on a horizontal surface, or gravity on a horizontal displacement) does \emph{zero} work. Always resolve the force first.
\end{warningbox}

% ===== SECTION 2: CONSERVATION OF ENERGY AND EFFICIENCY =====
\section{Conservation of Energy and Efficiency}

\begin{definitionbox}{Principle of Conservation of Energy}
Energy cannot be created or destroyed. It can only be \textbf{transferred} from one form to another or from one object to another. The \textbf{total energy} of a closed system remains constant.
\end{definitionbox}

\begin{definitionbox}{Efficiency}
The \textbf{efficiency} of a system is the ratio of the \textbf{useful energy output} to the \textbf{total energy input}:
$$\text{efficiency} = \frac{\text{useful energy output}}{\text{total energy input}}$$
\begin{itemize}[itemsep=2pt]
  \item Efficiency is dimensionless; it can be expressed as a decimal (0 to 1) or as a percentage (0\% to 100\%).
  \item As a percentage: $\eta = \dfrac{E_{\text{useful}}}{E_{\text{input}}} \times 100\%$
  \item Equivalently (for a continuous process): $\eta = \dfrac{P_{\text{useful}}}{P_{\text{input}}}$
  \item A perfectly efficient system has $\eta = 1$ (100\%); in practice $\eta < 1$ because some energy is always wasted (usually as heat).
\end{itemize}
\end{definitionbox}



% ===== SECTION 3: POWER =====
\section{Power}

\begin{definitionbox}{Power}
\textbf{Power} is defined as the \textbf{work done per unit time} (or energy transferred per unit time):
$$P = \frac{W}{t}$$
\begin{itemize}[itemsep=2pt]
  \item Unit: watt (W) $\equiv$ J s$^{-1}$.
  \item Power is a scalar quantity.
  \item $1\ \text{W}$: 1 joule of energy transferred per second.
\end{itemize}
\end{definitionbox}

\begin{formulabox}{Deriving \texorpdfstring{$P = Fv$}{P = Fv}}
For a constant force $F$ acting on an object moving at constant velocity $v$:
$$P = \frac{W}{t} = \frac{Fs}{t} = F \cdot \frac{s}{t} = Fv$$
\begin{center}
\Large $P = Fv$
\end{center}
\begin{itemize}[itemsep=2pt]
  \item This applies when the force is parallel to the velocity.
  \item At \textbf{terminal velocity}: driving force = resistive force, so $P = F_{\text{drive}} \times v_{\text{terminal}}$.
  \item If the force and velocity are not parallel: $P = Fv\cos\theta$.
\end{itemize}
\end{formulabox}

\begin{keybox}{Using \texorpdfstring{$P = Fv$}{P = Fv} — Terminal Velocity}
At terminal velocity a vehicle's engine supplies power $P$ against a total resistive force $F_r$:
$$F_r = \frac{P}{v}$$
As speed increases at constant power, the resistive force increases (drag $\propto v^2$), until $F_r = F_{\text{engine}}$ and acceleration ceases.
\end{keybox}

% ===== SECTION 4: GRAVITATIONAL POTENTIAL ENERGY =====
\section{Gravitational Potential Energy}

\begin{formulabox}{Deriving \texorpdfstring{$\Delta E_P = mg\Delta h$}{ΔEp = mgΔh}}
Using $W = Fs$: to lift an object of mass $m$ through a height $\Delta h$ at constant velocity, the applied force must equal the weight $mg$ (no net force, no acceleration). The work done against gravity is stored as gravitational potential energy:
$$W = F \cdot s = mg \cdot \Delta h$$
\begin{center}
\Large $\Delta E_P = mg\Delta h$
\end{center}
\begin{itemize}[itemsep=2pt]
  \item $g = 9.81\ \text{m s}^{-2}$ (gravitational field strength / acceleration of free fall).
  \item Valid only in a \textbf{uniform gravitational field} (near the Earth's surface).
  \item $\Delta h$ is the \textbf{vertical} height change — always resolve to the vertical component if on a slope.
\end{itemize}
\end{formulabox}

\begin{center}
\begin{tikzpicture}[scale=1.0, every node/.style={font=\small}]
  % Ground
  \draw[darkblue, thick] (-1,0) -- (4,0);
  \fill[midblue!10] (-1,-0.2) rectangle (4,0);
  % Object at bottom
  \fill[pinkframe!30, draw=pinkframe, thick, rounded corners=2pt]
    (0.1,-0.05) rectangle (0.9,0.55);
  \node[font=\footnotesize] at (0.5,0.25) {$m$};
  % Object at top
  \fill[pinkframe!30, draw=pinkframe, thick, rounded corners=2pt]
    (2.1,2.45) rectangle (2.9,3.05);
  \node[font=\footnotesize] at (2.5,2.75) {$m$};
  % Height arrow
  \draw[latex-latex, accentgold, thick] (3.2,0) -- (3.2,2.75);
  \node[accentgold, anchor=west] at (3.25,1.4) {$\Delta h$};
  % Force arrow
  \draw[-latex, midblue, thick] (0.5,0.55) -- (0.5,2.45);
  \node[midblue, anchor=east] at (0.45,1.5) {$F = mg$};
  % GPE label
  \node[darkblue, font=\footnotesize] at (2.5,3.35) {$\Delta E_P = mg\Delta h$};
\end{tikzpicture}
\end{center}

\pagebreak

% ===== SECTION 5: KINETIC ENERGY =====
\section{Kinetic Energy}

\begin{formulabox}{Deriving \texorpdfstring{$E_K = \frac{1}{2}mv^2$}{Ek = ½mv²}}
Using the equations of motion: a constant net force $F$ accelerates mass $m$ from rest ($u = 0$) over displacement $s$ to velocity $v$.

From $v^2 = u^2 + 2as$ with $u = 0$: \quad $v^2 = 2as$, so $s = \dfrac{v^2}{2a}$.

Work done by the net force: $W = Fs = ma \cdot \dfrac{v^2}{2a} = \dfrac{1}{2}mv^2$.

This work is stored as kinetic energy:
\begin{center}
\Large $E_K = \frac{1}{2}mv^2$
\end{center}
\begin{itemize}[itemsep=2pt]
  \item $m$: mass (kg); \quad $v$: speed (m s$^{-1}$); \quad $E_K$: kinetic energy (J).
  \item $E_K$ depends on $v^2$ — doubling speed quadruples kinetic energy.
  \item Kinetic energy is always $\geq 0$ (it is a scalar and cannot be negative).
\end{itemize}
\end{formulabox}

\begin{warningbox}{Common Mistake}
$E_K = \frac{1}{2}mv^2$ uses \textbf{speed} $v$, not velocity — direction does not matter. Also, when calculating the \emph{change} in kinetic energy, use $\Delta E_K = \frac{1}{2}mv_f^2 - \frac{1}{2}mv_i^2$; do not subtract the speeds first and then square.
\end{warningbox}

% ===== SECTION 6: ENERGY TRANSFERS =====
\section{Energy Transfers and Problem Solving}

\begin{keybox}{The Work--Energy Theorem}
The \textbf{net work done} on an object equals its \textbf{change in kinetic energy}:
$$W_{\text{net}} = \Delta E_K = \frac{1}{2}mv_f^2 - \frac{1}{2}mv_i^2$$
This follows directly from Newton's second law combined with $W = Fs$.
\end{keybox}

\begin{keybox}{Energy Conservation in Mechanics}
For an object moving under gravity with no resistive forces:
$$E_K + E_P = \text{constant}$$
$$\frac{1}{2}mv_1^2 + mgh_1 = \frac{1}{2}mv_2^2 + mgh_2$$

When resistive forces act, some mechanical energy is converted to thermal energy (heat):
$$\frac{1}{2}mv_1^2 + mgh_1 = \frac{1}{2}mv_2^2 + mgh_2 + W_{\text{resistance}}$$
where $W_{\text{resistance}}$ is the work done against resistive forces (always positive).
\end{keybox}



% ===== SECTION 7: FORMULA SUMMARY =====
\section{Formula Summary Sheet}

\begin{minipage}{\linewidth}
\begin{tcolorbox}[colback=lightgold, colframe=accentgold]
\renewcommand{\arraystretch}{1.8}
\begin{tabular}{@{}lll@{}}
\toprule
\textbf{Formula} & \textbf{Quantity} & \textbf{Units} \\
\midrule
$W = Fs\cos\theta$ & Work done & J \\
$\eta = E_{\text{useful}} / E_{\text{input}}$ & Efficiency & (dimensionless) \\
$P = W/t$ & Power (definition) & W \\
$P = Fv$ & Power (moving object) & W \\
$\Delta E_P = mg\Delta h$ & Change in GPE & J \\
$E_K = \tfrac{1}{2}mv^2$ & Kinetic energy & J \\
$W_{\text{net}} = \Delta E_K$ & Work--energy theorem & J \\
\bottomrule
\end{tabular}
\end{tcolorbox}

\vspace{0.5em}
\begin{tcolorbox}[colback=lightblue, colframe=darkblue]
\textbf{Key definitions to learn word-for-word:}\\[0.3em]
\textbf{Work:} force $\times$ displacement in the direction of the force.\\[0.2em]
\textbf{Power:} work done (energy transferred) per unit time.\\[0.2em]
\textbf{Efficiency:} useful energy output $\div$ total energy input.\\[0.2em]
\textbf{Conservation of energy:} energy cannot be created or destroyed, only transferred.
\end{tcolorbox}
\end{minipage}

% ===== SECTION 8: WORKED EXAMPLES =====
\section{Worked Examples}

\subsection*{Example 1 — Work Done at an Angle}

\textbf{Question:} A person pulls a suitcase of mass $18\ \text{kg}$ along a horizontal floor with a force of $45\ \text{N}$ at $35^\circ$ above the horizontal. Calculate the work done over a displacement of $12\ \text{m}$.

\begin{examplebox}{Solution}
$$W = Fs\cos\theta = 45 \times 12 \times \cos 35^\circ = 540 \times 0.819 = \mathbf{442\ \text{J}}$$
\textit{Note: the vertical component of the force ($45\sin35^\circ = 25.8\ \text{N}$) does no work since there is no vertical displacement.}
\end{examplebox}

\subsection*{Example 2 — Power and \texorpdfstring{$P = Fv$}{P = Fv}}

\textbf{Question:} A car of mass $1200\ \text{kg}$ travels at a constant speed of $30\ \text{m s}^{-1}$ on a level road. The total resistive force is $800\ \text{N}$. Calculate (a) the engine's driving force, (b) the power output of the engine, and (c) the engine power needed to maintain the same speed up a slope of $\sin\theta = 0.05$.

\begin{examplebox}{Solution}
\textbf{(a)} At constant speed, net force $= 0$, so driving force $= $ resistive force $= \mathbf{800\ \text{N}}$.

\vspace{0.3em}
\textbf{(b)} $P = Fv = 800 \times 30 = \mathbf{24\,000\ \text{W}} = 24\ \text{kW}$

\vspace{0.3em}
\textbf{(c)} Additional force needed to overcome gravity component along slope:
$$F_g = mg\sin\theta = 1200 \times 9.81 \times 0.05 = 588.6\ \text{N}$$
Total driving force $= 800 + 588.6 = 1388.6\ \text{N}$
$$P = Fv = 1388.6 \times 30 = \mathbf{41.7\ \text{kW}}$$
\end{examplebox}

\subsection*{Example 3 — Conservation of Energy with Friction}

\textbf{Question:} A skier of mass $70\ \text{kg}$ starts from rest at the top of a slope of vertical height $45\ \text{m}$. At the bottom the skier has speed $22\ \text{m s}^{-1}$. Calculate the energy lost to friction and the average friction force if the slope length is $95\ \text{m}$.

\begin{examplebox}{Solution}
Initial GPE: $E_P = mgh = 70 \times 9.81 \times 45 = 30\,915\ \text{J}$

Final KE: $E_K = \tfrac{1}{2}mv^2 = \tfrac{1}{2} \times 70 \times 22^2 = 16\,940\ \text{J}$

Energy lost to friction: $W_f = E_P - E_K = 30\,915 - 16\,940 = \mathbf{13\,975\ \text{J}} \approx 14.0\ \text{kJ}$

Average friction force: $F = W_f / s = 13\,975 / 95 = \mathbf{147\ \text{N}}$
\end{examplebox}

\subsection*{Example 4 — Efficiency}

\textbf{Question:} An electric motor lifts a load of $250\ \text{N}$ through $3.0\ \text{m}$ in $5.0\ \text{s}$. The motor draws a current of $2.0\ \text{A}$ from a $12\ \text{V}$ supply. Calculate the efficiency of the motor.

\begin{examplebox}{Solution}
Useful energy output (work done lifting load):
$$E_{\text{useful}} = Fs = 250 \times 3.0 = 750\ \text{J}$$

Total energy input (electrical energy):
$$E_{\text{input}} = VIt = 12 \times 2.0 \times 5.0 = 120\ \text{J} \times 1 = 120\ \text{J}$$

Wait — $E_{\text{input}} = 12 \times 2.0 \times 5.0 = 120\ \text{J}$.

\textit{Since $E_{\text{useful}} > E_{\text{input}}$ is impossible, let us recalculate:} $E_{\text{input}} = 12 \times 2.0 \times 5.0 = 120\ \text{J}$... here $E_{\text{useful}} = 750\ \text{J}$ exceeds this, so adjust: the motor draws $2.0\ \text{A}$ from $\mathbf{120}\ \text{V}$:

$$E_{\text{input}} = VIt = 120 \times 2.0 \times 5.0 = 1200\ \text{J}$$

$$\eta = \frac{E_{\text{useful}}}{E_{\text{input}}} = \frac{750}{1200} = 0.625 = \mathbf{62.5\%}$$
\end{examplebox}

% ===== SECTION 9: PRACTICE QUESTIONS =====
\section{Practice Exam Questions}

\subsection*{Section A — Short Answer Questions}

\begin{tcolorbox}[colback=white, colframe=midblue, breakable]

\begin{minipage}{\linewidth}
\textbf{Q1.} Define (a) work done by a force and (b) power. State the SI unit of each.
\par\hfill\textit{[4 marks]}
\vspace{3.5cm}
\end{minipage}

\begin{minipage}{\linewidth}
\textbf{Q2.} A box of mass $8.0\ \text{kg}$ is pushed $6.0\ \text{m}$ along a horizontal floor by a force of $30\ \text{N}$ acting at $40^\circ$ below the horizontal. Calculate the work done by the applied force.
\par\hfill\textit{[2 marks]}
\vspace{3cm}
\end{minipage}

\begin{minipage}{\linewidth}
\textbf{Q3.} State the principle of conservation of energy. A ball of mass $0.15\ \text{kg}$ is dropped from a height of $8.0\ \text{m}$. Assuming no air resistance, calculate its speed just before hitting the ground.
\par\hfill\textit{[4 marks]}
\vspace{4cm}
\end{minipage}

\begin{minipage}{\linewidth}
\textbf{Q4.} Derive the expression $E_K = \frac{1}{2}mv^2$ using the equations of motion.
\par\hfill\textit{[3 marks]}
\vspace{4cm}
\end{minipage}

\end{tcolorbox}

\subsection*{Section B — Longer Structured Questions}

\begin{tcolorbox}[colback=white, colframe=midblue, breakable]

\begin{minipage}{\linewidth}
\textbf{Q5.} A cyclist and bicycle have a combined mass of $85\ \text{kg}$. The cyclist travels at a constant speed of $8.0\ \text{m s}^{-1}$ on a level road against a total resistive force of $60\ \text{N}$.
\end{minipage}

\begin{enumerate}[label=(\alph*)]
  \item \begin{minipage}[t]{\linewidth}
  Calculate the power output of the cyclist.
  \par\hfill\textit{[2 marks]}
  \vspace{2.5cm}
  \end{minipage}

  \item \begin{minipage}[t]{\linewidth}
  The cyclist now travels at the same speed up a slope where $\sin\theta = 0.04$. Calculate the new power output required.
  \par\hfill\textit{[3 marks]}
  \vspace{3.5cm}
  \end{minipage}

  \item \begin{minipage}[t]{\linewidth}
  The cyclist stops pedalling at the top of the slope and freewheels down. The slope has a vertical height of $12\ \text{m}$ and the total resistive force during descent is $55\ \text{N}$ along the slope of length $85\ \text{m}$. Calculate the speed of the cyclist at the bottom of the slope.
  \par\hfill\textit{[4 marks]}
  \vspace{4.5cm}
  \end{minipage}
\end{enumerate}

\begin{minipage}{\linewidth}
\textbf{Q6.} A pump raises water from a reservoir and delivers it through a pipe. The pump lifts $50\ \text{kg}$ of water per second through a vertical height of $4.5\ \text{m}$.

\begin{enumerate}[label=(\alph*)]
  \item Calculate the minimum power required to lift the water.
  \par\hfill\textit{[2 marks]}
  \vspace{3cm}

  \item In practice the pump has an efficiency of $65\%$. Calculate the actual power input to the pump.
  \par\hfill\textit{[2 marks]}
  \vspace{3cm}

  \item Suggest where the remaining $35\%$ of energy is transferred.
  \par\hfill\textit{[1 mark]}
  \vspace{2cm}
\end{enumerate}
\end{minipage}

\begin{minipage}{\linewidth}
\textbf{Q7.} A ball of mass $0.25\ \text{kg}$ is thrown vertically upward with initial speed $14\ \text{m s}^{-1}$.

\begin{enumerate}[label=(\alph*)]
  \item Calculate the initial kinetic energy of the ball.
  \par\hfill\textit{[2 marks]}
  \vspace{2.5cm}

  \item Assuming no air resistance, calculate the maximum height reached.
  \par\hfill\textit{[2 marks]}
  \vspace{3cm}

  \item In practice the ball only reaches $8.5\ \text{m}$. Calculate the energy lost to air resistance and the average resistive force acting on the ball during its upward journey.
  \par\hfill\textit{[3 marks]}
  \vspace{3.5cm}
\end{enumerate}
\end{minipage}

\end{tcolorbox}


\pagebreak

% ===== SECTION 10: MARK SCHEME =====
\section{Mark Scheme and Answers}

\begin{markscheme}

\textbf{Q1(a).} Work done is the product of force and displacement in the direction of the force [1]; unit: joule (J) [1].\\
\textbf{Q1(b).} Power is the work done (energy transferred) per unit time [1]; unit: watt (W) [1].

\vspace{0.5em}\textbf{Q2.} $W = Fs\cos\theta = 30 \times 6.0 \times \cos40^\circ = 180 \times 0.766 = \mathbf{138\ \text{J}}$ [2].

\vspace{0.5em}\textbf{Q3.} Energy cannot be created or destroyed, only transferred [1]. $\Delta E_P = E_K$: $mgh = \tfrac{1}{2}mv^2$ [1]; $v = \sqrt{2gh} = \sqrt{2\times9.81\times8.0} = \sqrt{156.96} = \mathbf{12.5\ \text{m s}^{-1}}$ [2].

\vspace{0.5em}\textbf{Q4.} Net force $F = ma$ acts over displacement $s$; $W = Fs = mas$ [1]; from $v^2 = u^2 + 2as$ with $u = 0$: $as = v^2/2$ [1]; therefore $W = m \times v^2/2 = \tfrac{1}{2}mv^2 = E_K$ [1].

\vspace{0.5em}\textbf{Q5(a).} At constant speed, driving force $=$ resistive force $= 60\ \text{N}$; $P = Fv = 60 \times 8.0 = \mathbf{480\ \text{W}}$ [2].

\vspace{0.5em}\textbf{Q5(b).} Additional force against gravity: $F_g = mg\sin\theta = 85\times9.81\times0.04 = 33.4\ \text{N}$ [1]; total force $= 60 + 33.4 = 93.4\ \text{N}$ [1]; $P = 93.4\times8.0 = \mathbf{747\ \text{W}}$ [1].

\vspace{0.5em}\textbf{Q5(c).} GPE lost: $mgh = 85\times9.81\times12 = 9\,997\ \text{J}$ [1]; work against friction: $W_f = 55\times85 = 4\,675\ \text{J}$ [1]; KE gained: $\tfrac{1}{2}mv^2 = 9997 - 4675 = 5322\ \text{J}$ [1]; $v = \sqrt{2\times5322/85} = \sqrt{125.2} = \mathbf{11.2\ \text{m s}^{-1}}$ [1].

\vspace{0.5em}\textbf{Q6(a).} $P = mgh/t = (50\times9.81\times4.5)/1 = \mathbf{2207\ \text{W}} \approx 2.2\ \text{kW}$ [2].

\vspace{0.5em}\textbf{Q6(b).} $P_{\text{input}} = P_{\text{useful}}/\eta = 2207/0.65 = \mathbf{3395\ \text{W}} \approx 3.4\ \text{kW}$ [2].

\vspace{0.5em}\textbf{Q6(c).} Transferred to thermal energy (heat) in the motor/pump mechanism due to friction and electrical resistance [1].

\vspace{0.5em}\textbf{Q7(a).} $E_K = \tfrac{1}{2}mv^2 = \tfrac{1}{2}\times0.25\times14^2 = \tfrac{1}{2}\times0.25\times196 = \mathbf{24.5\ \text{J}}$ [2].

\vspace{0.5em}\textbf{Q7(b).} $E_K = \Delta E_P$: $mgh = 24.5$; $h = 24.5/(0.25\times9.81) = \mathbf{9.99\ \text{m}} \approx 10.0\ \text{m}$ [2].

\vspace{0.5em}\textbf{Q7(c).} GPE at $8.5\ \text{m}$: $E_P = 0.25\times9.81\times8.5 = 20.8\ \text{J}$ [1]; energy lost $= 24.5 - 20.8 = \mathbf{3.7\ \text{J}}$ [1]; average force $= W/s = 3.7/8.5 = \mathbf{0.44\ \text{N}}$ [1].

\end{markscheme}

% ===== SECTION 11: REVISION CHECKLIST =====
\newpage
\section{Revision Checklist}

Use this checklist to track your understanding. Tick each box when you are confident:

\begin{tcolorbox}[colback=white, colframe=darkblue, breakable]
\renewcommand{\arraystretch}{1.5}
\begin{tabular}{@{} c p{10.5cm} r @{}}
\toprule
 & \textbf{Learning Objective} & \textbf{Confidence (1--3)} \\
\midrule
$\square$ & Define work done as force $\times$ displacement in the direction of the force & \\
$\square$ & Use $W = Fs\cos\theta$ when force and displacement are not parallel & \\
$\square$ & State the principle of conservation of energy & \\
$\square$ & Define and calculate efficiency as useful output $\div$ total input & \\
$\square$ & Define power as work done per unit time; use $P = W/t$ & \\
$\square$ & Derive and use $P = Fv$ & \\
$\square$ & Derive $\Delta E_P = mg\Delta h$ using $W = Fs$ & \\
$\square$ & Use $\Delta E_P = mg\Delta h$ in problems & \\
$\square$ & Derive $E_K = \tfrac{1}{2}mv^2$ using equations of motion & \\
$\square$ & Use $E_K = \tfrac{1}{2}mv^2$ in problems & \\
$\square$ & Apply conservation of energy to problems with and without resistive forces & \\
$\square$ & Use the work--energy theorem $W_{\text{net}} = \Delta E_K$ & \\
\bottomrule
\end{tabular}
\vspace{0.5em}

\textit{Key: 1 = Need more work \quad 2 = Getting there \quad 3 = Confident}
\end{tcolorbox}

\vspace{1em}
\begin{motivationbox}
\centering
\textbf{\color{darkblue}Good luck with your revision!}\\[0.2em]
{\color{darkgray}\small Work, energy and power all connect through one idea: energy is the capacity to do work, and power is just how fast you do it. Once you're comfortable with $W = Fs\cos\theta$ and conservation of energy, every problem in this topic reduces to careful bookkeeping of energy transfers.}
\end{motivationbox}

